\honest{Burdensome Details}{Eliezer Yudkowsky}{2007}

I open with Pooh-Bah in \textsc{The Mikado}: ``Merely corroborative detail, intended to give artistic verisimilitude to an otherwise bald and unconvincing narrative.'' \nb{Tversky and Kahneman quote the same line in their 1983 paper, which the post cites.}

The conjunction fallacy is rating ``A and B'' as likelier than ``A'' alone, though a conjunction is never likelier than its parts. In one experiment, 68\% of subjects thought Reagan likelier to support unwed mothers and cut support to local governments than to support unwed mothers. The cause is that people ``substitute judgment of representativeness for judgment of probability,'' and a long series of experiments ``confirmed'' this. An added detail can make an outcome look more typical of what produces it: the unlikely claim and the likely one ``average out.'' So detail makes a story more plausible and less probable. We should then expect futurists to tell very plausible, detailed histories of the future, and people to accept large packages of unsupported claims wrapped around a few strong ones. \nb{That is Tversky and Kahneman's conclusion too: such additions ``can only lower probability.''}

Correcting yourself when you see both statements side by side is only a band-aid. In a 1982 experiment, professional forecasters gave a higher probability to ``Russia invades Poland, followed by suspension of diplomatic relations between the USA and the USSR'' than to the suspension alone. Each group saw only one statement, so no one could compare them. What could they have done? It seems to me that they would need to notice the word ``and'' and be ``shocked'' by it, and cut the probability by a factor of four at least, as the results imply. They might also have thought of all the reasons relations could be suspended, since the question was about any reason, not this one.

The Reagan subjects should add absurdities, measured in bits (the logarithm of the probability), rather than average them: cutting support to local governments may or may not happen (1 bit), supporting unwed mothers is very unlikely (4 bits), total 5 bits. Or simply: one strike and it is out. \nb{The bits add like this only if the two claims are independent. In a chain of causes, where each step makes the next likelier, adding them overstates the penalty.} In a 1983 experiment, subjects bet on which of three sequences would appear in twenty rolls of a die with four green faces and two red: RGRRR, GRGRRR or GRRRRR. 65\% chose GRGRRR, though any run containing it also contains RGRRR. They thought the sequence with the most green was best; they should have thought: ``Aha! Sequence 1 is short!'' \nb{Shortness wins here only because the longer sequence contains the shorter. With this die, six greens in a row is about 21 times likelier than five reds.} In general, you must ``feel every added detail as a burden.'' I give no evidence that feeling this way reduces the error.

Then a story. I once told a man taken with a futurist that universes replicating ``for any reason'' is likelier than their replicating through black holes that advanced civilizations make because universes evolve to make them do it. He said ``Oh.'' He had heard evidence that universes replicate and taken it as support for the whole package told around it. So hold up each detail and ask ``How do we know this detail?'' If someone describes a future nanotechnological war in which China breaks an international agreement, ask how they know it will be China.

The post ends, ``For it is written,'' with two lines: lighten your burden if you can, since any straw can break your back. \nb{They are from Yudkowsky's own ``Twelve Virtues of Rationality.''}
